\emph{Degree-three genus-one covers, without a Galois assumption.}
The same generated category also contains every degree-three
cover with genus-one source. We give the finite monodromy
argument explicitly. It is the genus-one case of the
classical trigonal construction described in
\citet[Section~9.1, Remark~9.13]{FF_RecillasRodriguez2003};
the finite verification below makes this consequence
independent of a convention about the genus in a
moduli-space statement.

Work over \(\overline{\mathbb Q}\). A connected
degree-three cover has geometric monodromy \(C_3\)
or \(S_3\). The first case is already abelian.
In the second case let \(a\) be the number of branch
transpositions and \(b\) the number of branch
three-cycles. Riemann--Hurwitz and the product-one
relation give
\[
          a+2b=6,\qquad a>0\text{ even},\qquad
          (a,b)\in\{(6,0),(4,1),(2,2)\}.
                                                        \tag{FF.48.12}
\]
Let \(q:S_4\twoheadrightarrow S_3\) be the action
on the three partitions of four letters into two
unordered pairs. Its kernel is the Klein four group.
Every branch tuple in (FF.48.12) that generates
\(S_3\) and has product one has a lift through
\(q\) that generates \(S_4\), has product one,
and preserves each permutation index. Here the index
of a permutation on \(d\) letters is \(d\) minus
its number of cycles.

This last assertion is a finite lemma verified by
complete enumeration in the supplied program. The
five nonidentity elements of \(S_3\) give
\(5^4+5^5+5^6=19375\) tuples of the possible
lengths. Filtering by index sum six, product one,
and generation of \(S_3\) leaves respectively
72, 270, and 240 tuples. A transposition has two
index-preserving lifts through \(q\), and a
three-cycle has four. For every one of the 582
retained tuples the program enumerates all 64
possible index-preserving lifts and finds exactly
twelve with product one and generated group
\(S_4\). It verifies the quotient map by all
\(24^2\) group products and saves an explicit
lift for every retained tuple. Thus the finite
assertion checks all cases, not just selected
ramification patterns. The program and certificates
are supplied at the following paths.

\path{agent_F/F_metabelian_checks.py}

\path{agent_F/F_recillas_certificates.json}

Apply the characteristic-zero Riemann existence
theorem to a chosen lift at the actual branch
points of the original cover. Its degree-four
source has
\(2g-2=-2\cdot4+6=-2\), hence genus zero.
Its Galois closure has group \(S_4\). The quotient
through \(q\) has exactly the original
\(S_3\) branch tuple, so the original cubic
cover is the subcover associated with the inverse
image of a point stabilizer in \(S_3\).
It therefore belongs to the category already
realized above.

Here descent from the complex construction retains the
covering diagram, not merely its branch data. Let
\(X=\mathbb P^1_{\overline{\mathbb Q}}\setminus B\),
where \(B\) is the actual finite branch set, enlarged
by the original affine infinity if necessary. The
Riemann existence theorem identifies its complex finite
etale covers with finite topological covering spaces
\citep[Expose~XII, Theoreme~5.1]{FF_SGA1}.
Base change from \(\overline{\mathbb Q}\) to
\(\mathbb C\) is an equivalence on finite etale
covering categories: apply
\citet[Expose~XIII, Proposition~4.6(a)]{FF_SGA1}
with \(k=\overline{\mathbb Q}\), the smooth curve
\(X\), and \(Y=\operatorname{Spec}\mathbb C\).
The required resolution hypothesis holds for curves,
and in characteristic zero the prime-to-characteristic
quotient is the entire fundamental group. Since
\(\pi_1(Y)=1\), the resulting fundamental-group map
is an isomorphism. This equivalence is fully faithful,
so the \(S_4\)-action, quotient diagram and isomorphism
with the original cubic cover all descend to
\(\overline{\mathbb Q}\). They use finitely many
algebraic coefficients and morphisms, and therefore
descend to one finite extension \(E'/E\).
Uniqueness of the scheme realization with the specified
actual point function then descends that realization
to \(E\), as detailed below.
